% --- IMPORTS ---
\documentclass[leqno]{article}
\usepackage{graphicx}
\usepackage{dsfont}
\usepackage{mathtools}
\usepackage{amssymb}
\usepackage{amsmath}
\usepackage{amsthm}
\usepackage{float}
\usepackage{ragged2e}
\usepackage{tensor}
\usepackage{fancyhdr}
\usepackage{empheq}
\usepackage[most]{tcolorbox}
\usepackage{enumitem}
\usepackage{dirtytalk}
\usepackage{marginnote}
\usepackage{microtype}
\usepackage[
  a4paper,
  left=0.8in,
  right=1in,
  top=1in,
  bottom=0.5in,
  headheight=14pt,
  headsep=0.4in
]{geometry}
\setlength{\parindent}{0cm}
% --- TikZ Packages ---
\usepackage{pgfplots}
\pgfplotsset{compat=1.18}
\usepgfplotslibrary{fillbetween, polar}
\usetikzlibrary{calc, positioning}
\usetikzlibrary{angles, quotes}
\usetikzlibrary{arrows.meta}
\usetikzlibrary{decorations.pathreplacing, decorations.markings}
\usetikzlibrary{patterns}
\usetikzlibrary{intersections}
% --- URL Packages ---
\usepackage[colorlinks=true,linkcolor=red,citecolor=magenta,urlcolor=black]{hyperref}%
\urlstyle{same}
% --- END OF IMPORTS ---

% --- CUSTOM COMMANDS ---
\RenewDocumentCommand{\marks}{o m}{%
  \IfValueTF{#1}%
    {\marginnote{\raggedright\textbf{[#2]}}[#1]}%
    {\marginnote{\raggedright\textbf{[#2]}}}%
}
\newcommand{\vspacerule}[2]{%
  \par\vspace*{#1}%
  \noindent\hrulefill\par
  \vspace*{#2}%
}
\definecolor{scarlet}{RGB}{205, 40, 75}
\newcommand{\questionitem}{%
  \item\phantomsection
  \addcontentsline{toc}{section}{Question \number\value{enumi}}%
}
\renewcommand{\contentsname}{Questions}
% --- END OF CUSTOM COMMANDS ---

% --- HEADER CONFIG ---
\pagestyle{fancy}
\fancyhead{}
\fancyfoot{}
\renewcommand{\headrulewidth}{0.9pt}
\lhead{Method of Differences}
\chead{}
\rhead{\href{https://danieljsmith.org}{danieljsmith.org}}
% --- END OF HEADER CONFIG ---

\begin{document}
\vspace*{20pt}
\tableofcontents
\newpage
\begin{enumerate}[
  label=\textbf{\arabic*.},
  leftmargin=*,
  topsep=0pt,
  partopsep=0pt,
  parsep=0pt
]

% ---- Question 1 ----
\questionitem

\begin{enumerate}[label=\textbf{(\alph*)}]
  \item Show that, for $r>0$,
  \[
  \ln\left(1+\frac{2}{r(r+3)}\right)=\ln\left(\frac{r+1}{r}\right)-\ln\left(\frac{r+3}{r+2}\right)
  \]
  \marks[-30pt]{1}
  \item Hence, using the method of differences, show that for integer $n\geq 1$,
  \[
  \sum_{r=1}^{n}\ln\left(1+\frac{2}{r(r+3)}\right)=\ln\left(\frac{3(n+1)}{n+3}\right)
  \]
  \marks[-30pt]{4}
\end{enumerate}

\newpage

% ---- Question 2 ----
\questionitem

Let $a$ be a positive constant. \\
\begin{enumerate}[label=\textbf{(\alph*)}]
  \item Use the method of differences to find
  \[
  \sum_{r=1}^{n} \frac{1}{(r+a)(r+a+1)(r+a+2)}
  \]
  in terms of $n$ and $a$. \marks{4} \\
  \item Hence find the value of $a$ for which
  \[
  \sum_{r=1}^{\infty} \frac{1}{(r+a)(r+a+1)(r+a+2)} = \frac{1}{24}
  \]
  \marks[-30pt]{3}
\end{enumerate}

\newpage

% ---- Question 3 ----
\questionitem

\begin{enumerate}[label=\textbf{(\alph*)}]
  \item Show that, for $r \ge 2$,
  \[
  \frac{4r}{\sqrt{r(r+2)}+\sqrt{r(r-2)}}=\sqrt{r(r+2)}-\sqrt{r(r-2)}
  \]
  \marks[-30pt]{2}
  \item Hence use the method of differences to determine
  \[
  \sum_{r=2}^{n} \frac{4r}{\sqrt{r(r+2)}+\sqrt{r(r-2)}}
  \]
  giving your answer in simplest form. \marks{3} \\
  \item Hence show that
  \[
  \sum_{r=3}^{7} \frac{4r}{\sqrt{r(r+2)}+\sqrt{r(r-2)}} = A\sqrt{2}+B\sqrt{3}+C\sqrt{7}
  \]
  where $A$, $B$ and $C$ are integers to be determined. \marks{2} \\
\end{enumerate}

\newpage

% ---- Question 4 ----
\questionitem

Use the method of differences to show that
\[
\sum_{r=1}^{n} \frac{12}{(2r-1)(2r+1)(2r+3)} = \frac{n(an+b)}{(2n+1)(2n+3)}
\]
where $a$ and $b$ are integers to be determined. \marks{6} \\

\newpage

% ---- Question 5 ----
\questionitem

\begin{enumerate}[label=\textbf{(\alph*)}]
  \item Use the method of differences to prove that for $n \ge 1$
  \[
  \sum_{r=1}^{n} \ln\left(1+\frac{2}{r(r+3)}\right)=\ln\left(\frac{3(n+1)}{n+3}\right)
  \]
  \marks[-30pt]{4}
  \item Hence find the exact value of
  \[
  \sum_{r=25}^{74} \ln\left(\left(1+\frac{2}{r(r+3)}\right)^{28}\right)
  \] \\
  Give your answer in the form $a\ln\left(\dfrac{b}{c}\right)$ where $a$, $b$ and $c$ are integers to be determined. \marks{3} \\
\end{enumerate}

\newpage

% ---- Question 6 ----
\questionitem

Using partial fractions and a method of differences, prove that for positive integer $n$,
\[
\sum_{r=1}^{n} \frac{6}{r(r+1)(r+2)} = \frac{n(an+b)}{2(n+1)(n+2)}
\]
where $a$ and $b$ are constants to be found. \marks{6} \\

\newpage

% ---- Question 7 ----
\questionitem

Let $T_n = (n+1)^3 + (n+2)^3 + \cdots + (2n)^3$. \\
\begin{enumerate}[label=\textbf{(\alph*)}]
  \item By considering $(r+1)^4-(r-1)^4$, use the method of differences to prove that
  \[
  \sum_{r=1}^{n} r^3 = \frac{1}{4}n^2(n+1)^2
  \]
  \marks[-30pt]{5}
  \item Show that \[T_n = \frac{1}{4}n^2(an^2+bn+c)\] where $a$, $b$ and $c$ are integers to be determined. \marks{3} \\
  \item State the value of \[\lim_{n \to \infty} \frac{T_n}{n^4}\] \marks{1} \\
\end{enumerate}

\newpage

% ---- Question 8 ----
\questionitem

\begin{enumerate}[label=\textbf{(\alph*)}]
  \item By expressing $\dfrac{1}{r(r+2)}$ in partial fractions, show that
  \[
  \sum_{r=1}^{n} \frac{1}{r(r+2)} = \frac{3}{4} - \frac{1}{2(n+1)} - \frac{1}{2(n+2)}
  \]
  \marks[-30pt]{5}
  \item Hence determine the value of
  \[
  \sum_{r=1}^{\infty} \frac{1}{r(r+2)}
  \]
  \marks[-20pt]{2}
\end{enumerate}

\newpage

% ---- Question 9 ----
\questionitem

Show that
\[
\frac{1}{1 \times 4} + \frac{1}{2 \times 5} + \frac{1}{3 \times 6} + \cdots + \frac{1}{n(n+3)} = \frac{11}{18} - \frac{an^2+bn+c}{3(n+1)(n+2)(n+3)}
\]
where $a$, $b$ and $c$ are constants to be found. \marks{6} \\

\newpage

% ---- Question 10 ----
\questionitem

\begin{enumerate}[label=\textbf{(\alph*)}]
  \item Use standard results from the formula booklet to show that
  \[
  \sum_{r=1}^{N} r(2r-1)(2r+1)=\frac{1}{2}N(N+1)(2N^2+2N-1)
  \]
  \marks[-30pt]{3}
  \item Express $\dfrac{6r+13}{(2r-1)(2r+3)}$ in partial fractions and hence use the method of differences to find
  \[
  \sum_{r=1}^{N} \frac{6r+13}{(2r-1)(2r+3)}\left(\frac{1}{2}\right)^{r+1}
  \]
  in terms of $N$. \marks{4} \\
  \item Deduce the value of
  \[
  \sum_{r=1}^{\infty} \frac{6r+13}{(2r-1)(2r+3)}\left(\frac{1}{2}\right)^{r+1}
  \]
  \marks[-20pt]{1}
\end{enumerate}

\newpage

% ---- Question 11 ----
\questionitem

\begin{enumerate}[label=\textbf{(\alph*)}]
  \item Show that
  \[
  (2n)^5-(2n-2)^5 \equiv 10(2n-1)^4+20(2n-1)^2+2
  \]
  \marks[-20pt]{2}
  \item Hence, using the method of differences, show that for all positive integer values of $n$,
  \[
  \sum_{r=1}^{n}(2r-1)^4=\frac{1}{15}n(2n-1)(2n+1)(an^2+bn+c)
  \]
  where $a$, $b$ and $c$ are integers to be determined. \marks{7} \\
\end{enumerate}

\newpage

% ---- Question 12 ----
\questionitem

Let
\[
S_n = \sum_{r=1}^{n} \frac{1}{(3r-2)(3r+4)}
\]
where $n \geq 1$. \\

\begin{enumerate}[label=\textbf{(\alph*)}]
  \item Use the method of differences to show that
  \[
  S_n = \frac{n(15n+17)}{8(3n+1)(3n+4)}
  \]
  \marks[-30pt]{6}
  \item Show that, for any number $k$ greater than $\dfrac{24}{5}$, if the difference between $\dfrac{5}{24}$ and $S_n$ is less than $\dfrac{1}{k}$, then \\[5pt]
  \[
  n > \frac{k - 15 + \sqrt{k^2 + 81}}{18}
  \]
  \marks[-30pt]{6}
\end{enumerate}
\end{enumerate}
\end{document}
